\section{Restricted Mertens sums and two-dimensional Abel summation}
\label{sec:double-abel}

\subsection{Zero-extended differences}

For a finitely supported kernel \(F:\N^2\to\mathbb C\), set
\begin{align*}
 \Delta_1F(x,y)&=F(x,y)-F(x+1,y),\\
 \Delta_2F(x,y)&=F(x,y)-F(x,y+1),\\
 \Delta_{12}F(x,y)
 &=F(x,y)-F(x+1,y)-F(x,y+1)+F(x+1,y+1).
\end{align*}
All values outside the declared support are zero.  Define the rectangular
mixed variation
\begin{equation}
 \cV_\square(F)
 =
 \sum_{x,y\ge1}|\Delta_{12}F(x,y)|.
\label{eq:mixed-variation}
\end{equation}
Because of zero extension, \(\cV_\square\) includes all upper and right
boundary jumps and the terminal corner.  No separate edge term is omitted.

\begin{lemma}[Two-dimensional Abel identity]
\label{lem:double-abel}
Let \((\alpha_x)\) and \((\beta_y)\) be finitely supported sequences, and
put
\[
 A(x)=\sum_{r\le x}\alpha_r,
 \qquad
 B(y)=\sum_{s\le y}\beta_s.
\]
For every finitely supported \(F\),
\begin{equation}
 \boxed{
 \sum_{x,y\ge1}\alpha_x\beta_yF(x,y)
 =
 \sum_{x,y\ge1}A(x)B(y)\Delta_{12}F(x,y).}
\label{eq:double-abel}
\end{equation}
\end{lemma}

\begin{proof}
One-dimensional summation by parts with zero extension gives
\[
 \sum_x\alpha_xG(x)
 =
 \sum_xA(x)\bigl(G(x)-G(x+1)\bigr).
\]
Apply this first in \(x\) for each \(y\), and then in \(y\).  The resulting
second difference is exactly \(\Delta_{12}F\).
\end{proof}

\begin{remark}[Equivalent edge--corner form]
If a kernel is specified only on a closed rectangle and not zero-extended,
expanding \eqref{eq:double-abel} produces an interior mixed-difference sum,
two terminal edge sums, and one corner term.  Definition
\eqref{eq:mixed-variation} packages all four pieces without losing them.
\end{remark}

\subsection{Restricted Mertens envelopes}

For \(d\ge1\), define
\begin{equation}
 M_d(x)
 =
 \sum_{\substack{n\le x\\(n,d)=1}}\mu(n),
 \qquad
 \mathfrak M_d(X)
 =
 \max_{1\le x\le X}|M_d(x)|.
\label{eq:restricted-mertens}
\end{equation}
We set \(\mathfrak M_d(X)=0\) when \(X<1\).
These quantities depend only on the squarefree kernel of \(d\):
\[
 M_d=M_{\operatorname{rad}(d)}.
\]
We use only the classical prime-number-theorem consequence
\[
 M_1(x)=o(x);
\]
see, for example, \citet[Chapter 5]{IwaniecKowalski2004}.

\begin{lemma}[Fixed local factors preserve Mertens cancellation]
\label{lem:fixed-d-mertens}
For every fixed positive integer \(d\),
\begin{equation}
 M_d(x)=o_d(x),
 \qquad
 \mathfrak M_d(X)=o_d(X).
\label{eq:fixed-d-mertens}
\end{equation}
\end{lemma}

\begin{proof}
It suffices to adjoin one prime at a time.  If \(p\nmid d\), partition the
sum defining \(M_d(x)\) according to divisibility by \(p\).  For the
\(p\)-divisible terms, write \(n=pr\).  A nonzero term requires
\((r,dp)=1\), and then \(\mu(pr)=-\mu(r)\).  Therefore
\[
 M_{dp}(x)=M_d(x)+M_{dp}(x/p).
\]
Iteration gives
\begin{equation}
 M_{dp}(x)=\sum_{k\ge0}M_d(x/p^k),
\label{eq:restricted-mertens-recursion}
\end{equation}
where the sum is finite.

If \(M_d(t)=o(t)\), split \eqref{eq:restricted-mertens-recursion} at a
fixed large threshold.  Above the threshold, the geometric sum is
\(o(x)\); below it, the trivial estimate \(|M_d(t)|\le t\) gives a bounded
geometric tail, also \(o(x)\).  Starting with \(d=1\) proves the first
claim by induction over the prime divisors of \(d\).

For the maximal claim, given \(\varepsilon>0\), take \(x_0\) such that
\(|M_d(x)|\le\varepsilon x\) for \(x\ge x_0\).  Then
\[
 \mathfrak M_d(X)
 \le
 \max_{x<x_0}|M_d(x)|+\varepsilon X
\]
Divide by \(X\), take the limit superior, and then let
\(\varepsilon\to0\).  This proves \(\mathfrak M_d(X)=o(X)\).
\end{proof}

\begin{remark}[No growing-\(d\) uniformity]
\Cref{lem:fixed-d-mertens} holds for each fixed \(d\).  It does not state
a uniform estimate when \(d\) grows with the physical scale.  The master
certificate below avoids silently assuming such uniformity by truncating at
fixed \(D\) and retaining the remainder as an explicit tail.
\end{remark}

\begin{proposition}[Abel form of one diagonal slice]
\label{prop:one-diagonal-abel}
For
\[
 K^{[d]}(x,y)=K(dx,dy),
\]
one has
\begin{align}
 &\sum_{\substack{x,y\ge1\\(x,d)=(y,d)=1}}
 \mu(x)\mu(y)K(dx,dy)
\notag\\
 &\hspace{2cm}=
 \sum_{x,y\ge1}
 M_d(x)M_d(y)\Delta_{12}K^{[d]}(x,y).
\label{eq:one-diagonal-abel}
\end{align}
Consequently,
\begin{equation}
 \left|
 \sum_{\substack{x,y\ge1\\(x,d)=(y,d)=1}}
 \mu(x)\mu(y)K(dx,dy)
 \right|
 \le
 \mathfrak M_d(U/d)\mathfrak M_d(V/d)
 \cV_\square(K^{[d]})
\label{eq:one-diagonal-bound}
\end{equation}
when \(K\) is supported on \([1,U]\times[1,V]\).
\end{proposition}

\begin{proof}
Apply \cref{lem:double-abel} to
\[
 \alpha_x=\mu(x)\one_{(x,d)=1},
 \qquad
 \beta_y=\mu(y)\one_{(y,d)=1}.
\]
Although these sequences are infinite, the kernel is supported on a finite
rectangle; truncate both sequences beyond that rectangle before applying
the lemma.
Their partial sums are \(M_d(x)\) and \(M_d(y)\).
\end{proof}
